\scanpage{8}
161.

Next, setting
\[
y'=y+(q/3\alpha)x^2-(u/3\alpha)x^3\qquad(u\text{ arbitrary})
\]
we find similarly that in (6) we may assume further that $q=ux$.

Finally, setting
\[
x'=x+vy
\]
and choosing suitable values for $u,v$ we find that in (6) we can take $p=q=r=s=t=0$, proving (ii).

(ii) $\Rightarrow$ (iii). Since $R^*/m^*=R/m$ is algebraically closed of characteristic $\ne3,5$, and $R^*$ is henselian, $\alpha$ is a cube in $R^*$ and $\beta$ is a fifth power. Since $m^*=mR^*$ it is now clear that (ii) $\Rightarrow$ (iii).

(iii) $\Rightarrow$ (iv) $\Rightarrow$ (i). Since $\widehat R$ is also the completion of $R^*$, it follows at once from (iii), in view of the Cohen structure theorem, that $\widehat R$ is a homomorphic image of $\overline R=S/(u^2+v^3+w^5)$. But in CASE Vd of \S24 we have seen that $\overline R$ is factorial; since $\dim.\widehat R=\dim.\overline R=2$, we must therefore have $\widehat R=\overline R$.\qquad q.e.d.

REMARK 25.2. If $R/m$ has characteristic 2, 3, or 5, we must change Theorem 25.1 somewhat. (i) and (i)' remain the same, but for the relation in (ii) we have two or more possibilities, as indicated below, with $\alpha$ and $\beta$ units in $R$. (iii) will now state that $\alpha$ and $\beta$ can both be assumed to be 1 if $R$ is henselian. The corresponding change in (iv) is obvious.
\mspage{7}
